
\documentclass[12pt,a4paper]{article}

%------------------------------------------------
% PAQUETES
%------------------------------------------------
\usepackage[spanish]{babel}
\usepackage[utf8]{inputenc}
\usepackage[T1]{fontenc}
\usepackage{geometry}
\usepackage{amsmath}
\usepackage{amssymb}
\usepackage{booktabs}
\usepackage{array}
\usepackage{enumitem}
\usepackage{xcolor}

\geometry{margin=2.5cm}
\setlength{\parskip}{0.5em}
\setlength{\parindent}{0pt}

%------------------------------------------------
% DOCUMENTO
%------------------------------------------------
\begin{document}

\begin{titlepage}
    \centering

    \vspace*{3cm}

    {\Large Universidad de Buenos Aires\par}
    \vspace{0.3cm}
    {\large Facultad de Ciencias Exactas y Naturales\par}

    \vspace{2cm}

    {\huge\bfseries TP2: AirMiles\par}
    \vspace{0.5cm}
    {\LARGE Grupo JavaConMates\par}

    \vspace{2cm}

    \textbf{Integrantes:}\\[0.3cm]
    Ayala Ignacio\\
    Della Bonzana Lara\\
    Morales Santiago\\
    Vannelli Tiziana

    \vfill

    {\large \today\par}
\end{titlepage}

%------------------------------------------------
\section{Introducción}
%------------------------------------------------

En este documento se justifican las estructuras de datos
seleccionadas para implementar el TAD \texttt{AirMiles}.
Para cada estructura se describe la información que almacena,
las operaciones que permite realizar y su relación con las
complejidades temporales esperadas.

Denotamos por $U$ la cantidad de usuarios registrados,
por $P$ la cantidad de promociones y por $|L|$ la longitud
del nombre de una promoción. Además, $N$ representa la
cantidad de usuarios solicitados para el ranking.

El diseño busca que las operaciones de consulta y modificación
accedan directamente a la información que necesitan, evitando
recorridos innecesarios. Las complejidades que dependen del
ABB se analizan distinguiendo el caso promedio del peor caso,
ya que el árbol utilizado no realiza balanceo automático.

%------------------------------------------------
\section{Almacenamiento e identificación de usuarios}
%------------------------------------------------

\subsection{Estructura elegida}

Los usuarios se almacenan en un árbol binario de búsqueda
(ABB). Cada nodo contiene un objeto \texttt{Usuario}, y el
orden del árbol se determina por el identificador único
(\texttt{id}).

La clase \texttt{Usuario} implementa
\texttt{Comparable<Usuario>}. Su método \texttt{compareTo}
compara los identificadores mediante
\texttt{Integer.compare}. Esto permite que el ABB localice
un usuario a partir de su ID sin confundirlo con otro usuario
que tenga distintos valores de millas.

\subsection{Justificación}

El ABB permite organizar los usuarios por identificador y
realizar búsquedas e inserciones siguiendo un único camino
desde la raíz. En un árbol razonablemente equilibrado, estas
operaciones requieren tiempo $O(\log U)$.


\subsection{Operaciones asociadas}

El ABB se utiliza para:

\begin{itemize}[nosep]
    \item \texttt{registrarUsuario} y
    \texttt{registrarUsuarioConReferidor}: insertar usuarios.
    \item \texttt{consultarSaldo} y
    \texttt{consultarCategoria}: localizar al usuario por ID.
    \item \texttt{acumularMillas}: encontrar la cuenta que se
    debe actualizar.
    \item \texttt{canjearMillas}: localizar al usuario que
    realiza el canje.
    \item \texttt{transferirMillas}: localizar al usuario
    de origen y al de destino.
    \item \texttt{paresTransferidoresExclusivos}: permitiría
    recorrer todos los usuarios, si la implementación de esta
    operación lo requiere.
\end{itemize}

La operación auxiliar \texttt{buscarUsuario} centraliza el
acceso al ABB y evita repetir la lógica de búsqueda en los
distintos métodos de \texttt{AirMiles}.

%------------------------------------------------
\section{Representación de cada usuario}
%------------------------------------------------

Cada usuario se representa mediante un objeto
\texttt{Usuario}. El objeto agrupa los datos de su cuenta,
por lo que, una vez localizado en el ABB, las operaciones
pueden consultar o modificar sus atributos directamente.

Entre los atributos considerados se encuentran:

\begin{itemize}[nosep]
    \item \texttt{id}: identificador del usuario.
    \item \texttt{millasDisponibles}: saldo que puede utilizar.
    \item \texttt{millasTotales}: millas acumuladas históricamente.
    \item \texttt{millasCanjeadas}: millas utilizadas en canjes.
    \item \texttt{millasTransferidas}: millas transferidas a
    otros usuarios.
    \item \texttt{posicionHeap}: posición actual del usuario
    en el max-heap.
\end{itemize}

Los atributos que representan cantidades históricas permiten
actualizar las estadísticas sin tener que reconstruirlas a
partir de todas las operaciones anteriores. Por su parte,
\texttt{posicionHeap} permite localizar al usuario dentro
del heap en tiempo $O(1)$ antes de actualizar su prioridad.

La comparación de usuarios por ID se mantiene independiente
de su saldo: el ABB necesita un orden estable basado en la
identidad del usuario, mientras que el heap utiliza el saldo
disponible para determinar las prioridades.

%------------------------------------------------
\section{Promociones}
%------------------------------------------------

\subsection{Estructura elegida}

Las promociones se almacenan en un trie o árbol de prefijos.
Cada nodo contiene un arreglo de 26 hijos, uno por cada letra
del alfabeto, y puede almacenar la información de una promoción
cuando representa el final de su nombre.

La información de cada promoción se agrupa en un objeto
\texttt{Promocion}, que contiene:

\begin{itemize}[nosep]
    \item \texttt{factor}: multiplicador de la promoción.
    \item \texttt{tope}: cantidad máxima de millas que se pueden
    otorgar utilizando la promoción, cuando tiene límite.
    \item \texttt{millasOtorgadas}: cantidad de millas que ya
    consumieron el tope.
    \item \texttt{tieneTope}: indica si existe un límite.
\end{itemize}

La promoción inicial \texttt{NombrePromo}, con factor $1$,
se agrega al crear el sistema.

\subsection{Justificación}

El trie permite localizar una promoción recorriendo los
caracteres de su nombre. Si el acceso a los hijos de cada
nodo cuesta $O(1)$, buscar, insertar y eliminar una promoción
requiere $O(|L|)$, independientemente de la cantidad total
de promociones almacenadas.

Esta estructura es adecuada porque las operaciones sobre
promociones dependen de sus nombres. Además, agrupar los datos
de cada promoción en un mismo objeto permite consultar el
factor y controlar el tope sin mantener estructuras separadas
para cada atributo.

\subsection{Control del tope}

Para acumular millas se calcula primero la cantidad base,
aplicando el factor de categoría del usuario. Luego, el método
\texttt{calcularMillas} busca la promoción una sola vez y
calcula cuántas millas se obtendrían al aplicar su factor.

Si la promoción no tiene tope, se aplica el factor normalmente.
Si tiene tope, se verifica que la cantidad acumulada más las
millas de la nueva operación no supere el límite. Cuando la
operación es válida, se actualiza \texttt{millasOtorgadas}.
Si superaría el tope, se otorgan únicamente las millas base y
el acumulado de la promoción no se modifica.

Por lo tanto, el tope cuenta todas las millas otorgadas en
las acumulaciones en las que se aplica la promoción, no solo
las millas adicionales producidas por el multiplicador.

Este control requiere una única búsqueda en el trie, por lo
que la complejidad de esta parte es $O(|L|)$.

La implementación supone que los nombres de promoción están
formados por letras de la A a la Z, sin espacios ni otros
caracteres. Si la especificación admite un alfabeto más amplio,
la representación de los hijos debe adaptarse.

%------------------------------------------------
\section{Máximos históricos}
%------------------------------------------------

Las operaciones \texttt{elMasAcumulador} y
\texttt{elMasCanjeador} deben consultar usuarios que alcanzan
determinados máximos históricos.

Para evitar recorrer el ABB en cada consulta, se mantienen
referencias directas a los usuarios que ocupan esas posiciones:
\texttt{maxAcumulador} y \texttt{maxCanjeador}.

Cuando un usuario acumula millas, se actualiza su contador
histórico y se compara con el máximo registrado. De manera
análoga, cuando un usuario canjea o transfiere millas, se
actualiza la información que corresponda al criterio definido
por la especificación para el máximo canjeador.

Si el usuario supera el máximo actual, basta con reemplazar
la referencia correspondiente. La consulta posterior consiste
en acceder a esa referencia y devolver el ID, por lo que cuesta
$O(1)$.

Al registrar el primer usuario, las referencias deben
inicializarse para evitar consultar un máximo que todavía sea
\texttt{null}. También debe respetarse el criterio de desempate
que establezca la especificación.

Esta estrategia evita recorrer todos los usuarios en cada
consulta, a cambio de mantener las referencias actualizadas
durante las operaciones que modifican los contadores.

%------------------------------------------------
\section{Ranking de millas disponibles}
%------------------------------------------------

\subsection{Estructura elegida}

Para obtener los usuarios con mayor saldo disponible se
utiliza un max-heap. La prioridad de cada elemento se determina
por \texttt{millasDisponibles}; el usuario con mayor saldo se
encuentra en la raíz.

El heap es genérico y recibe las operaciones necesarias para
comparar elementos y consultar o actualizar sus posiciones.
De esta forma, puede funcionar con objetos \texttt{Usuario}
sin depender de la implementación de las cuentas.

\subsection{Justificación}

La propiedad de max-heap garantiza que la prioridad de cada
padre es mayor o igual que la de sus hijos. No exige que el
hijo izquierdo sea mayor que el derecho ni mantiene a todos
los elementos ordenados.

Esta propiedad permite consultar el máximo en $O(1)$ y extraerlo
en $O(\log U)$. Para obtener los primeros $N$ usuarios, se
extrae repetidamente el máximo. Como las extracciones modifican
temporalmente el heap, los elementos extraídos se guardan y
luego se vuelven a insertar, de modo que el ranking no altera
el conjunto de usuarios registrados.

Tanto las extracciones como las reinserciones cuestan
$O(\log U)$ cada una. Por eso, obtener y restaurar los primeros
$N$ elementos cuesta $O(N\log U)$, suponiendo $N\leq U$.

\subsection{Actualización de prioridades}

Cuando cambia el saldo disponible de un usuario, su prioridad
en el heap también puede cambiar. Para actualizarla
eficientemente, cada usuario conserva su posición actual
mediante \texttt{posicionHeap}.

Así se evita buscar el usuario recorriendo todo el heap.
Una vez conocida su posición, se restaura la propiedad de
max-heap moviendo el elemento hacia arriba o hacia abajo
según corresponda, en tiempo $O(\log U)$.

La posición se actualiza cada vez que se intercambian elementos.
Cuando un usuario se extrae, su posición se marca como
inválida; al reinsertarlo, se establece nuevamente.

La complejidad $O(N\log U)$ supone que el heap contiene a
todos los usuarios y que se mantiene actualizado después de
cada operación que modifica el saldo.

%------------------------------------------------
\section{Transferencias entre usuarios}
%------------------------------------------------

La operación \texttt{transferirMillas} localiza al usuario de
origen y al de destino mediante el ABB. Luego modifica los
saldos disponibles de ambas cuentas y actualiza los contadores
históricos que correspondan.

Si el ranking se mantiene en un max-heap, ambas prioridades
deben actualizarse después de modificar los saldos. Cada
búsqueda en el ABB cuesta $O(\log U)$ en promedio y $O(U)$
en el peor caso; cada actualización del heap cuesta
$O(\log U)$.

Para identificar pares de transferidores exclusivos no basta
con conocer los totales transferidos por cada usuario. Es
necesario conservar información sobre las relaciones entre
origen y destino, de modo que se pueda distinguir entre
transferencias recíprocas, transferencias hacia terceros y
repeticiones entre una misma pareja.

La estructura auxiliar elegida para estas relaciones debe
permitir verificar la exclusividad y evitar contar dos veces
un mismo par. La implementación de
\texttt{paresTransferidoresExclusivos} debe completar este
mantenimiento y demostrar que la consulta cumple la
complejidad exigida por la especificación.

%------------------------------------------------
\section{Complejidades temporales}
%------------------------------------------------

La siguiente tabla resume las complejidades derivadas de
las estructuras descriptas. Las operaciones que consultan o
modifican usuarios dependen de la altura del ABB; por eso se
indica tanto el caso promedio como el peor caso.

\begin{center}
\renewcommand{\arraystretch}{1.3}
\begin{tabular}{|p{7.0cm}|p{5.4cm}|}
\hline
\textbf{Operación} & \textbf{Complejidad temporal} \\
\hline
Crear el sistema & $O(1)$ para inicializar las estructuras \\
\hline
Registrar usuario & $O(\log U)$ promedio; $O(U)$ peor caso \\
\hline
Buscar usuario & $O(\log U)$ promedio; $O(U)$ peor caso \\
\hline
Acumular millas & $O(\log U+|L|)$ promedio; $O(U+|L|)$ peor caso \\
\hline
Canjear millas & $O(\log U)$ promedio; $O(U)$ peor caso \\
\hline
Transferir millas & $O(\log U)$ promedio; $O(U)$ peor caso \\
\hline
Consultar saldo o categoría & $O(\log U)$ promedio; $O(U)$ peor caso \\
\hline
Iniciar o finalizar promoción & $O(|L|)$ \\
\hline
Controlar tope de promoción & $O(|L|)$ \\
\hline
Consultar máximo histórico & $O(1)$ \\
\hline
Consultar máximo del heap & $O(1)$ \\
\hline
Actualizar prioridad en el heap & $O(\log U)$ \\
\hline
Obtener y restaurar los primeros $N$ del ranking & $O(N\log U)$ \\
\hline
Pares de transferidores exclusivos & Depende de completar la estructura auxiliar \\
\hline
\end{tabular}
\end{center}

Las complejidades de las operaciones que utilizan el ABB
suponen que las implementaciones de sus métodos básicos
cumplen las cotas indicadas. Como el árbol no se balancea
automáticamente, no se garantiza $O(\log U)$ en el peor caso.

Las complejidades de las operaciones del heap suponen que
las comparaciones y las consultas o actualizaciones de
posición cuestan $O(1)$.

%------------------------------------------------
\section{Conclusión}
%------------------------------------------------

Las estructuras elegidas responden a necesidades diferentes
del TAD. El ABB organiza los usuarios por identificador y
permite localizar sus cuentas; el trie permite administrar
promociones en función de la longitud de sus nombres; y el
max-heap permite obtener los usuarios con mayor saldo sin
ordenar todas las cuentas en cada consulta.

Las referencias a los máximos históricos permiten responder
esas consultas en tiempo constante, mientras que la posición
almacenada por cada usuario facilita mantener el heap
actualizado cuando cambian los saldos.

La corrección del sistema depende de mantener consistentes
estas estructuras auxiliares durante las operaciones. En
particular, deben actualizarse los saldos, los contadores
históricos, las prioridades del heap y el consumo de los topes
promocionales. La representación de las relaciones de
transferencia debe completarse de acuerdo con la condición
de exclusividad establecida por la especificación.

\end{document}
```
